% ============================================================
% tab/tabla_comparativa.tex
%
% COLUMNAS OBLIGATORIAS. No sustituir, no reordenar, no añadir.
% Se usa table* porque con siete columnas no cabe en una sola
% columna del formato IEEE.
%
% Cada fila sale de los campos 4 a 7 de la plantilla de extracción:
% unidad de análisis, protocolo, métrica y resultado reportado.
%
% NO incluir aquí referencias puramente metodológicas (calibración,
% leakage genérico, imputación): no resuelven una tarea comparable
% sobre un dataset y obligarían a inventar valores. Esas viven en
% el texto.
% ============================================================
\begin{table*}[!t]
\centering
\caption{Resumen comparativo de trabajos relacionados.}
\label{tab:comparativa}
\footnotesize
\setlength{\tabcolsep}{4pt}
\begin{tabularx}{\textwidth}{@{}l c P{2.4cm} P{2.4cm} P{2.0cm} P{1.8cm} Y@{}}
\toprule
\textbf{Referencia} & \textbf{Año} & \textbf{Técnica} &
\textbf{Dataset / entorno} & \textbf{Métrica principal} &
\textbf{Resultado reportado} & \textbf{Limitación identificada}\\
\midrule
\pc{cita} & \pc{} & \pc{} & \pc{} & \pc{} & \pc{} & \pc{}\\
\pc{cita} & \pc{} & \pc{} & \pc{} & \pc{} & \pc{} & \pc{}\\
\pc{cita} & \pc{} & \pc{} & \pc{} & \pc{} & \pc{} & \pc{}\\
\pc{cita} & \pc{} & \pc{} & \pc{} & \pc{} & \pc{} & \pc{}\\
\pc{cita} & \pc{} & \pc{} & \pc{} & \pc{} & \pc{} & \pc{}\\
\bottomrule
\end{tabularx}
\end{table*}